% Candidate theorem and comparison section. No priority or advantage claim.
\section{Implicit completion of single-recurrent-boundary components}

\paragraph{Input and boundary.}
The input is an explicit finite e-graph with nonnegative binary-encoded integer
node weights and a set of required e-classes. A feasible extraction selects one
representative in each reached class and has no selected dependency cycle.
Costs count a reached representative once. Semantic equivalence is an input
assumption; mapped hardware performance is not the objective in this section.
Let $C$ contain the required classes and every root-reachable class with at
least two distinct possible parent classes. The remaining private dependency
regions are acyclic, and regions owned by distinct boundary classes are disjoint
until another boundary is reached. We use the earlier ownership lemma.

Project possible dependency paths through private regions onto $C$ and compute
the strongly connected components $B_1,\ldots,B_m$ of this support graph.
A local implementation of $v\in C$ expands private classes and stops at boundary
references. An implementation containing a reference to $v$ itself cannot
participate in an acyclic extraction and may be discarded.

\paragraph{Structural promise.}
Every remaining local implementation of every $v$ references at most one
\emph{distinct other} class in $B(v)$, its support component. It may reference
arbitrarily many boundary classes outside $B(v)$. Repeated ports to the same
boundary class count as one dependency. This promise includes every
acyclic-support input but also genuine cyclic alternatives. It does not assert
that the entire program is unary. The promise can be checked in polynomial time:
propagate local support labels $\emptyset$, a singleton in $B(v)$, or a sentinel
for two-or-more. Discard self-references, combine distinct children by union and
alternatives by choice. The private graph is acyclic. A surviving sentinel is
an obstruction, not an infeasibility certificate. The recognition may conservatively
reject globally unusable conjunctive alternatives.

\paragraph{Implicit local tables.}
Fix an active set $A\subseteq C$ containing all required classes. Core references
outside $A$ are forbidden. For each $v\in A$, compute $w_v(\bot;A)$ and
$w_v(u;A)$ for $u\in (A\cap B(v))\setminus\{v\}$. These are minimum private
costs of local implementations using, respectively, no same-component boundary
reference, or exactly the single reference $u$. References in other components
are permitted exactly when in $A$ and have zero local price. Their actual cost
will be charged at their own boundary class. The tables use min-plus propagation
with at most $|B(v)|+1$ numerical entries at a private vertex; impossible entries
are infinite. No list of all attainable costs, threshold grid, or cost spectrum
is constructed. Conditional minima and witnesses suffice.

\paragraph{Arborescence completion theorem.}
For each component $B$, make a directed graph on
$\{\bot\}\cup(A\cap B)$. Include $\bot\to v$ with weight $w_v(\bot;A)$
and $u\to v$ with weight $w_v(u;A)$ whenever finite. Let $\tau_B(A)$ be the
minimum weight of a spanning out-arborescence rooted at $\bot$, or infinity
if none exists. Define $F(A)=\sum_B\tau_B(A)$, with the empty component costing
zero. Under the structural promise, $F(A)$ equals the minimum cost of an
acyclic computation constructing all of $A$ and no other boundary classes.
Consequently,
\[
 \operatorname{OPT}=\min_{A\supseteq R}F(A),
\]
where $R$ is the required boundary set. Any sound mandatory-class inference may
replace $R$ by a larger set which every valid extraction contains.

\paragraph{Proof.}
In a local implementation there is at most one same-component prerequisite.
Reverse the dependency arrow to point from prerequisite to result. Selecting
one local implementation per active boundary class therefore selects one incoming
arc at every non-ground vertex. Acyclicity is exactly the condition that these
arcs form a rooted spanning arborescence: every predecessor chain terminates at
$\bot$ rather than a cycle. Replacing any local implementation by the cheapest
one with its same internal prerequisite does not change that condition. External
prerequisites are in strictly earlier components of the dependency condensation
and cannot create an intercomponent cycle. Thus any computation for $A$ has
cost at least $\sum_B\tau_B(A)$. Conversely, optimal component arborescences
choose local witnesses. Construct them dependency-first within each arborescence
and dependency-first across the condensation. All selected prerequisites then
exist. Private ownership prevents double charging or conflicting representative
choices. This constructs all of $A$ at the stated cost. Taking the active boundary
of an optimum extraction proves one direction of the last equality. For the
other direction, remove unused pieces from a constructed $A$; nonnegative weights
ensure the root-reachable extraction costs no more.\hfill$\square$

\paragraph{Complexity and quantum corollary.}
Minimum arborescence is a standard polynomial-time problem (Chu--Liu/Edmonds;
Tarjan). A simple one-cycle-at-a-time implementation costs $O(VE)$ arithmetic
operations. The local tables and their witnesses are polynomial in the original
explicit graph, not a compiled numerical catalog. With $N$ the explicit size,
$L$ the cost bit length including a threshold $K$, and $h=|C\setminus R|$,
activation enumeration runs in $O(2^h\operatorname{poly}(N,L))$ time and
polynomial workspace. Standard bounded-error quantum search for $F(A)\leq K$
gives
\[
 O(2^{h/2}\operatorname{poly}(N,L))
\]
time and polynomial workspace in the \emph{original} input. Multiplying all
weights and $K$ by an integer changes arithmetic bit lengths, not the activation
domain. A uniform register for $h$ optional bits is directly preparable. The
polynomial classical predicate can be compiled into a fixed-time Boolean circuit
and made reversible by keeping and uncomputing its bounded history. Dynamic
indices can use explicit linear scans, with polynomial overhead, rather than
free QRAM. A measured activation is completed and checked classically. No
complete Toffoli counts or fault-tolerance resources are claimed here. This is
a speedup over activation enumeration, not a lower bound against all classical
extractors or a claimed new quantum-search theorem.

\paragraph{Why the promise matters.}
Let all $X,Y,Z$ be required. $X$ can be built for ten or from both $Y,Z$ for
zero; each of $Y,Z$ can be built for one or from $X$ for zero. The true optimum
is two. Replacing the two-prerequisite rule by a choice of one predecessor makes
an invalid cost-one arborescence appear. The implementation reports this input
as outside its promise; it does not silently drop or relax the conjunction.

\paragraph{Classical comparative boundary.}
If \emph{every original representative} has at most one distinct child, extraction
is a directed Steiner tree problem: use a ground edge for a leaf and a weighted
edge from the prerequisite to its resulting class, with required outputs as
terminals. Both directions preserve optimum cost. The standard terminal-subset
DP supplies a different reference upper bound
$O(n^3+3^r n+2^r n^2)$ for $r$ terminals, using $O(2^r n+n^2)$ space.
A single output is just a shortest-path problem, even with cyclic alternatives.
These classical reductions must not be suppressed when comparing against
$2^h$ enumeration. Treewidth/pathwidth extraction, simplification, heuristic
and exact native solvers remain additional competitors on the broader promise.
No useful-workload prevalence, priority for the combined reduction, or
best-classical separation has been established.
